\appendix
\section*{Appendix A: Proof of Proposition 3}

\paragraph{A.1}
This appendix provides the detailed derivation of the
comparative statics stated in Proposition~3
(\S\ref{eq:32:cs_omega}--\ref{eq:32:cs_price}). The proof
proceeds by differentiating the optimal mechanization gap
\eqref{eq:32:kappa_star} with respect to each state
variable.

\subsection*{A.1 The Optimal Mechanization Gap}

\paragraph{A.2}
Recall from \S3.2 that the firm maximizes the net
viability surplus:
\begin{equation}
	\max_{\kappa_t}\; c(\kappa_t)
	= \Phi(\kappa_t) - \kappa_t \cdot MC(\omega_{t-1}, p^\kappa_t)
	\label{eq:A:objective}
\end{equation}
where $\Phi(\kappa_t) = \lambda_0 + \lambda_1 \kappa_t
+ \lambda_2 \kappa_t^2$ is the mechanization frontier
\eqref{eq:32:frontier_kappa}, and
$MC(\omega_{t-1}, p^\kappa_t) = \gamma_0 - \gamma_1 \omega_{t-1}
+ \gamma_2 p^\kappa_t$ is the marginal viability cost
\eqref{eq:32:MC}. Substituting and differentiating:
\begin{align}
	c(\kappa_t) &= \lambda_0
	+ \left[\lambda_1 - \gamma_0 + \gamma_1 \omega_{t-1}
	- \gamma_2 p^\kappa_t\right] \kappa_t
	+ \lambda_2 \kappa_t^2
	\label{eq:A:unconstrained} \\
	\frac{dc}{d\kappa_t} &= \lambda_1 - \gamma_0
	+ \gamma_1 \omega_{t-1} - \gamma_2 p^\kappa_t
	+ 2\lambda_2 \kappa_t^* = 0
	\label{eq:A:foc}
\end{align}
Solving \eqref{eq:A:foc} under Assumption~1
($\lambda_2 < 0$) yields the optimal mechanization gap:
\begin{equation}
	\kappa_t^*(\omega_{t-1}, p^\kappa_t)
	= \frac{(\gamma_0 - \lambda_1) - \gamma_1 \omega_{t-1}
		+ \gamma_2 p^\kappa_t}{2\lambda_2}
	\label{eq:A:kappa_star}
\end{equation}
which reproduces \eqref{eq:32:kappa_star}.

\subsection*{A.2 The Distributive Channel}

\paragraph{A.3}
\textbf{Proof of} \eqref{eq:32:cs_omega}. Differentiating
\eqref{eq:A:kappa_star} with respect to $\omega_{t-1}$:
\begin{equation}
	\frac{\partial \kappa_t^*}{\partial \omega_{t-1}}
	= \frac{-\gamma_1}{2\lambda_2}
	\label{eq:A:cs_omega}
\end{equation}
Under Assumption~1, $\lambda_2 < 0$. By Assumption~3
($\gamma_1 > 0$), the numerator $-\gamma_1 < 0$ and the
denominator $2\lambda_2 < 0$. Therefore:
\begin{equation}
	\frac{\partial \kappa_t^*}{\partial \omega_{t-1}}
	= \frac{-\gamma_1}{2\lambda_2} > 0
	\label{eq:A:cs_omega_sign}
\end{equation}
A rising wage share forces cost-minimizing firms to widen
the mechanization gap, substituting machinery for labor.
$\hfill\blacksquare$

\subsection*{A.3 The Relative Price Channel}

\paragraph{A.4}
\textbf{Proof of} \eqref{eq:32:cs_price}. Differentiating
\eqref{eq:A:kappa_star} with respect to $p^\kappa_t$:
\begin{equation}
	\frac{\partial \kappa_t^*}{\partial p^\kappa_t}
	= \frac{\gamma_2}{2\lambda_2}
	\label{eq:A:cs_price}
\end{equation}
Under Assumption~1, $\lambda_2 < 0$. By Assumption~3
($\gamma_2 > 0$), the numerator $\gamma_2 > 0$ and the
denominator $2\lambda_2 < 0$. Therefore:
\begin{equation}
	\frac{\partial \kappa_t^*}{\partial p^\kappa_t}
	= \frac{\gamma_2}{2\lambda_2} < 0
	\label{eq:A:cs_price_sign}
\end{equation}
A rising relative price of machinery suppresses the optimal
mechanization gap, as the structural cost of the composition
shift penalizes mechanization. $\hfill\blacksquare$

\subsection*{A.4 The Cost Hurdle Formulation (Alternative
	Derivation)}

\paragraph{A.5}
An equivalent derivation expresses the FOC as a cost
hurdle. Define the effective cost hurdle of mechanization:
\begin{equation}
	H(\omega_{t-1}, p^\kappa_t)
	\equiv \frac{MC(\omega_{t-1}, p^\kappa_t)}{1}
	= \gamma_0 - \gamma_1 \omega_{t-1} + \gamma_2 p^\kappa_t
	\label{eq:A:hurdle}
\end{equation}
The FOC \eqref{eq:A:foc} can be rewritten as:
\begin{equation}
	\Phi'(\kappa_t^*) = H(\omega_{t-1}, p^\kappa_t)
	\label{eq:A:foc_hurdle}
\end{equation}
The firm mechanizes up to the point where the marginal
productivity of the mechanization gap equals the effective
cost hurdle. The comparative statics follow from the
implicit function theorem applied to
\eqref{eq:A:foc_hurdle}:
\begin{align}
	\frac{\partial \kappa_t^*}{\partial \omega_{t-1}}
	&= -\frac{\partial H / \partial \omega_{t-1}}
	{\Phi''(\kappa_t^*)}
	= -\frac{-\gamma_1}{2\lambda_2} > 0
	\label{eq:A:ift_omega} \\
	\frac{\partial \kappa_t^*}{\partial p^\kappa_t}
	&= -\frac{\partial H / \partial p^\kappa_t}
	{\Phi''(\kappa_t^*)}
	= -\frac{\gamma_2}{2\lambda_2} < 0
	\label{eq:A:ift_price}
\end{align}
Both derivations yield identical results, confirming
Proposition~3.